% Section 6.6, from lectures/L06/appendix/e_exercises.md.

\section{Exercises}
\label{c6:sec:exercises}

\begin{aside}
\textbf{How to check your work.} Exercise~\ref{c6:ex:driver} is checked by this chapter's test
suite: write the file at the path the specification gives, then run \code{make test}. The suite is
cumulative and runs Chapter~1 to Chapter~5's tests too. Exercise~\ref{c6:ex:capstone} is the C
capstone, which no test covers: it needs the C runtime and is built by hand, with the command in
\secref{c6:sec:asmfromc}.

The rest will have worked solutions in the companion repository, published later as
\file{lectures/L06/appendix/f_solutions.md}, and the hand calculations have short answers in
Appendix~\ref{app:answers}. Work each one before you read it.

\textbf{No hardware is needed for any of this.}
\end{aside}

Each exercise is labelled with its kind. There is exactly one \textbf{Cross-check}, and here it is
Exercise~\ref{c6:ex:crosscheck}. It is the last one in the book.

% ------------------------------------------------------------------------------------------
\exercise{What it catches}{Recall}
\label{c6:ex:catches}
\begin{parts}
  \item State, in one sentence, the only thing a watchdog detects.
  \item For each of the following, say whether a watchdog would catch it:
    \begin{itemize}
      \item a loop that never exits;
      \item a driver writing to the wrong port register;
      \item the torn 16-bit read from Chapter~3;
      \item a program waiting for an interrupt that will never arrive;
      \item a stack that has grown into the variables.
    \end{itemize}
  \item One of those five is ``sometimes, eventually''. Which, and what has to happen first?
  \item Give the two modes a timeout can produce, and say which of them is a safety mechanism.
\end{parts}

% ------------------------------------------------------------------------------------------
\exercise{Where to pet it}{Design}
\label{c6:ex:pet}
\begin{parts}
  \item A program has a main loop and a timer interrupt firing every millisecond. Somebody puts
    \code{watchdog_reset} in the timer handler, because it is called reliably. Explain precisely
    what that achieves and what it destroys.
  \item Where should it go instead, and what property does that location need?
  \item The main loop has two branches, one of which runs rarely. Where do you pet the watchdog so
    that a fault in the rare branch is still caught?
  \item A program pets the watchdog in three different places. Say what has probably gone wrong
    with its design.
\end{parts}

% ------------------------------------------------------------------------------------------
\exercise{Control bytes}{Hand calculation}
\label{c6:ex:bytes}
Using \secref{c6:sec:registers}, and without running anything:
\begin{parts}
  \item Give the byte for a 1 second timeout in system reset mode.
  \item Give the byte for an 8 second timeout in interrupt mode.
  \item Give the byte for a 250\,ms timeout in interrupt-and-reset mode.
  \item Somebody wants four seconds and writes \code{0x08}. Say exactly what they asked for
    instead, and describe what the device does from that moment on.
  \item Why is their mistake so much worse than, say, asking for 2 seconds and getting 250\,ms?
\end{parts}

\noindent\textbf{Check yourself:} the timeout table is in \secref{c6:sec:registers}, and the four
selector bits are not contiguous, so no control byte is ever the timeout number.

% ------------------------------------------------------------------------------------------
\exercise{The timed sequence}{Recall}
\label{c6:ex:sequence}
\begin{parts}
  \item Give the two writes, in order, and the constraint between them.
  \item Four cycles is the limit. Roughly how many cycles would an interrupt take if one landed
    between the two writes? Say why \code{cli} is therefore not optional.
  \item On a device fresh from reset, what is in \code{WDTCSR} if the second write arrives too
    late? Say why that particular value is the worst possible outcome.
  \item Why must \code{WDRF} be cleared before \code{WDE} can be? Give the reason the hardware
    behaves that way, and the consequence for a driver that forgets.
  \item The sequence begins with \code{wdr}. What would go wrong without it?
\end{parts}

% ------------------------------------------------------------------------------------------
\exercise{The watchdog driver}{Code}
\label{c6:ex:driver}
Write \file{drivers/source/watchdog.asm} to the specification in \secref{c6:sec:driver}.
\begin{parts}
  \item Write all four subroutines.
  \item Run \code{make test}.
  \item Insert four \code{nop} instructions between the two writes of the sequence, rebuild, and
    note what \code{WatchdogDriver.InitWritesTheByteAsked} reports. Say what that value means, then
    remove them.
  \item Replace the save-and-restore of \code{SREG} with a plain \code{cli} and \code{sei}. Note
    which test fails and describe the bug it is protecting you from.
\end{parts}

% ------------------------------------------------------------------------------------------
\exercise{Calling it from C}{Code}
\label{c6:ex:capstone}
Write \file{drivers/app/main.c} to the specification in \secref{c6:sec:capstone}.
\begin{parts}
  \item Write a C program that calls \code{shift_bits} and does something with the result. Change
    nothing in \file{utils.asm}; the \file{utils_gnu.S} you translate from it is a transcription,
    not a redesign.
  \item Translate \file{utils.asm} into \file{drivers/source/utils_gnu.S} using the table in
    \secref{c6:sec:asmfromc}, then build with \code{avr-gcc}, giving it both the \file{.c} and the
    \file{.S}. This build shares no flag with the book's usual one: say what
    \code{-mmcu=atmega328p} is deciding that \code{avra} decided some other way, and say what the
    startup code does that your \file{main.asm} used to do by hand.
  \item Declare the prototype's argument as \code{uint16_t} instead, rebuild, and say what now
    happens for an argument of 5 and for an argument of 300. Then declare the result as
    \code{uint16_t} as well and answer the same question. Then put both back.
  \item Write a C function and call it from your assembly. Say which registers you had to assume it
    would destroy.
\end{parts}

% ------------------------------------------------------------------------------------------
\exercise{Sleeping}{Design}
\label{c6:ex:sleep}
\begin{parts}
  \item For each of idle, power-down and power-save, say what is still running.
  \item A program sleeps in power-down and expects Timer1 to wake it after one second. Describe
    what actually happens, and how long it lasts.
  \item Name the two wake sources that survive every sleep mode, and say why each of them does.
  \item Design, in words, a device that does something once every eight seconds and spends the
    rest of its life drawing microamps. Say which peripheral wakes it and in which mode.
  \item A device refuses to stay asleep: it wakes immediately, every time. Give two plausible
    causes.
\end{parts}

% ------------------------------------------------------------------------------------------
\exercise[fillaccent]{Your assembly against the compiler's}{Cross-check}
\label{c6:ex:crosscheck}
\emph{Compute it by hand, measure it, reconcile.}

The last cross-check in the book, and the only one where the thing you are checking against is
somebody else's work.
\begin{parts}
  \item \textbf{By hand.} From Chapter~1, your \code{shift_bits} costs $6n + 10$ cycles
    (\secref{c1:sec:onpaper}). Write that down.
  \item \textbf{Write the same thing in C.} A function taking a byte, returning \code{1} shifted
    left that many times, using a loop. \textbf{Call it \code{c_shift_bits}}, so that both live in
    one ELF without colliding. Compile it at \code{-Os} and read the assembly with
    \code{avr-gcc -S}. Count the instructions and predict its cost as a formula in $n$, before
    measuring.
  \item \textbf{By measurement.} \file{drivers/build/drivers.hex} holds no C, so build one image
    that holds both, the same hand build the capstone uses, and point \code{make measure} at it:
\begin{shell}
avr-gcc -mmcu=atmega328p -Os \
        drivers/app/main.c drivers/source/utils_gnu.S \
        -o drivers/build/mixed.elf
make measure IMAGE=mixed \
     CALLS="shift_bits:0 c_shift_bits:0 shift_bits:8 c_shift_bits:8"
\end{shell}
    \code{IMAGE} names the stem and not the extension, so the same variable reaches
    \file{drivers.hex} and this \file{mixed.elf}; \file{ci/measure.sh} takes whichever of the two
    exists.
  \item \textbf{Reconcile.} Give both formulas. They are not the same shape, and one of them wins
    at some arguments and loses at others: say where they cross.
  \item \textbf{The surprise.} At one specific argument the compiler's version is \textbf{faster}
    than yours. Which argument, by how much, and what did the compiler do that you did not? Look
    at the instruction it used and say what it does.
  \item \textbf{Code size.} Count the instructions in each. Now say which version you would ship,
    and note that ``the hand-written one is smaller and faster'' is not the answer the numbers
    support.
  \item \textbf{What this is evidence for.} In one paragraph, say what this comparison does and
    does not establish about writing assembly by hand. Be specific about the size of the win and
    about what it cost you to get it.
\end{parts}

% ------------------------------------------------------------------------------------------
\exercise{What you would actually reach for}{Design}
\label{c6:ex:reach}
The book is over. For each of the following, say whether you would write assembly, write C, or do
something else entirely, and why.
\begin{parts}
  \item A driver for a new sensor on the SPI bus.
  \item An interrupt handler that must respond within two microseconds.
  \item A routine that toggles a pin for a fixed 3 microseconds, exactly, every time.
  \item The main loop of a battery-powered logger.
  \item Working out why a program that worked yesterday now resets every few seconds.
  \item One of the five answers above is the one this book was really for. Say which.
\end{parts}
